% Source: https://artofproblemsolving.com/community/c3h1683835_do_you_understand_this_help
% Stable identity: aops-topic-1683835-q01
Multiplying an equation by a nonzero number, or dividing it by one, gives an equivalent equation.

If $\dfrac {a}{b}=\dfrac {c}{d}$, multiply both sides by $bd\ne0$: $\dfrac {a}{b}\cdot bd=ad$ and $\dfrac {c}{d}\cdot bd=bc$, so $ad=bc$.

Conversely, if $ad=bc$, divide both sides by $bd\ne0$: $\dfrac{ad}{bd}=\dfrac {a}{b}$ and $\dfrac{bc}{bd}=\dfrac {c}{d}$, so $\dfrac {a}{b}=\dfrac {c}{d}$. $\blacksquare$

In other words, cross-multiplying just writes both fractions over the common denominator $bd$: $\dfrac {a}{b}=\dfrac{ad}{bd}$ and $\dfrac {c}{d}=\dfrac{bc}{bd}$, and two fractions with the same denominator are equal exactly when their numerators are equal.

For $\dfrac{5}{8}$ and $\dfrac{20}{32}$: $5\cdot32=160=8\cdot20$, so $\boxed{\dfrac{5}{8}=\dfrac{20}{32}}$.
